\label{sec:comparison}

\subsection{HH vs.\ Webster: Geometric vs.\ Arithmetic Mean}

The only difference between Huntington-Hill and Webster is the rounding
threshold: $\sqrt{s(s+1)}$ (HH) vs.\ $s + 0.5$ (Webster).
For all $s \geq 1$:
\[
  s < \sqrt{s(s+1)} < s + 0.5.
\]
The gap between the thresholds is:
\[
  (s+0.5) - \sqrt{s(s+1)}
  = (s+0.5) - \sqrt{s^2+s}
  = \frac{(s+0.5)^2 - (s^2+s)}{(s+0.5) + \sqrt{s^2+s}}
  = \frac{0.25}{(s+0.5) + \sqrt{s^2+s}}
  \approx \frac{1}{8s}
\]
for large $s$.
This gap is small for large $s$ (large states, where the threshold
rarely matters) and largest for small $s$ (small states).
For $s = 1$: HH threshold $= \sqrt{2} \approx 1.414$, Webster threshold
$= 1.5$; the gap is 0.086.

In practice, the two methods agree for all but a small number of states
near the boundary.
For the 2020 Census, Huntington-Hill and Webster agree on all 50 states
(J.2 identifies the specific states closest to the boundary);
the difference between the two methods manifests in some historical
censuses and for smaller House sizes.

\begin{table}[h]
\centering
\caption{Rounding thresholds for seats 1--6 under Huntington-Hill and Webster.}
\label{tab:thresholds}
\begin{tabular}{rrrr}
\toprule
Seats $s$ & HH threshold $\sqrt{s(s+1)}$ & Webster threshold $s+0.5$ & Gap \\
\midrule
1 & 1.4142 & 1.5000 & 0.0858 \\
2 & 2.4495 & 2.5000 & 0.0505 \\
3 & 3.4641 & 3.5000 & 0.0359 \\
4 & 4.4721 & 4.5000 & 0.0279 \\
5 & 5.4772 & 5.5000 & 0.0228 \\
6 & 6.4807 & 6.5000 & 0.0193 \\
\bottomrule
\end{tabular}
\end{table}

Because HH's threshold is lower than Webster's, a state whose exact quota
$q_i = p_i/d$ falls between $\sqrt{s(s+1)}$ and $s+0.5$ will receive
$s+1$ seats under HH but only $s$ seats under Webster.
States with small populations (and thus small $s$) are more likely
to fall in this interval, which is why HH is described as
slightly small-state-favoring relative to Webster.

\subsection{HH vs.\ Adams: Geometric Mean vs.\ Ceiling}

Adams uses ceiling rounding: state $i$ receives $\lceil p_i/d \rceil$
seats, meaning the threshold is at $s$ (i.e., any non-integer quota
rounds up to the next integer).
For small states, Adams is substantially more generous than HH
because Adams never rounds a fractional quota down.
For example, a state with exact quota 1.05 receives 2 seats under Adams
but only 1 seat under HH (since HH's threshold is $\sqrt{1\cdot 2}
\approx 1.414 > 1.05$, the state does not yet qualify for a second seat).

The practical effect on 2020 Census data: Alaska (quota $\approx 0.96$)
receives 1 seat under both Adams and HH (the constitutional floor provides
the 1-seat minimum regardless); Rhode Island (quota $\approx 1.44$)
receives 2 seats under Adams and 2 seats under HH;
Montana at 2020 (quota $\approx 1.42$) receives 2 seats under both.
Adams would diverge from HH for a state with quota between 1 and $\sqrt{2}
\approx 1.414$, which does not occur in the 2020 Census but occurred in
some historical censuses.

\subsection{Summary: Divisor Method Spectrum}

The four divisor methods lie on a spectrum ordered by their rounding
thresholds.
For a state with exact quota $q \in (s, s+1)$, the state receives $s+1$
seats if its quota exceeds the method's threshold:

\begin{center}
\begin{tabular}{lll}
\toprule
Method & Threshold for $(s+1)$-th seat & Bias \\
\midrule
Adams           & $q > s$ (always rounds up)     & Most small-state-favoring \\
Huntington-Hill & $q > \sqrt{s(s+1)}$            & Slightly small-state-favoring \\
Webster         & $q > s + 0.5$                  & Neutral \\
Jefferson       & $q > s + 1$ (always rounds down) & Most large-state-favoring \\
\bottomrule
\end{tabular}
\end{center}

Huntington-Hill lies between Webster and Adams on this spectrum,
reflecting Congress's choice of a method that is neither strongly
pro-small-state (Adams) nor neutral (Webster),
but slightly protective of smaller states' proportional representation.
